\scanpage{147}
\mspage{721}

even an $(\mathfrak S_3,\mathcal J^0)$-torsor (a $\mathcal J^0$-torsor on which $\mathfrak S_3$ operates in a way compatible with its operations on $\mathcal J^0$), hence classified by an element of $H^1(S_0,\mathfrak S_3;\mathcal J^0)$. Now I presume that we have
\[
\tag{307}
H^0(S_0,\mathfrak S_3,\mathcal J^0)=H^1(S_0,\mathfrak S_3,\mathcal J^0)=0
\]
(where
\[
\tag{308}
\left\{\begin{aligned}
S_0&=\operatorname{Spec}\mathbb Z,\quad \mathcal J^0=\mathbb G_m\times_{\mathbb Z}M\simeq\mathbb G_m^{\{0,1,\infty\}}/\mathrm{diag}\\
&\simeq\operatorname{Ker}\left(\diagram{\begin{tikzcd}[ampersand replacement=\&,column sep=small]
\mathbb G_m^{\{0,1,\infty\}}\arrow[r] \& \mathbb G_m^{\deleted{\unsure{\{0,1,\infty\}}}}\\
(x_0,x_1,x_\infty)\arrow[r,mapsto] \& \substack{\inserted{x_0\cdot x_1\cdot x_\infty}\\\deleted{\unsure{(x_1x_\infty^{-1},x_\infty x_0^{-1},x_0x_1^{-1})}}}
\end{tikzcd}}\right)
\end{aligned}\right.
\]
\margin{Careful, this last isomorphism is obtained \unsure{via}
\diagram{\begin{tikzcd}[ampersand replacement=\&]
(x_0,x_1,x_\infty)\arrow[r,mapsto] \& (x_1x_\infty^{-1},x_\infty x_0^{-1},x_0x_1^{-1})
\end{tikzcd}}, which commutes with $\mathfrak S_3^+$, but \unsure{not|not} with $\mathfrak S_3$ as a whole}

We have indeed
\[
\begin{aligned}
H^0(S_0,\mathfrak S_3;\mathcal J^0)&\simeq\left\{(x_0,x_1,x_\infty)\in\mathbb Z^{*\{0,1,\infty\}}\ \middle|\ \begin{array}{l}x_0x_1x_\infty=1\\(x_0,x_1,x_\infty)\text{ invariant under }\mathfrak S_3\end{array}\right\}\\
&=\{(x,x,x)\mid x\in\mathbb Z^{*\{0,1,\infty\}},\quad x^3=1\}\\
&=1
\end{aligned}
\]
\margin{Calculation \unsure{correct}, for it \unsure{rests on this} $\mathfrak S_3^+$-isomorphism which is not an $\mathfrak S_3$-isomorphism, \unsure{it seems here to be|there seems not to be} need of $H^0(S,\mu_3)=1$, whereas a calculation already made of the \unsure{splittings} of extensions (cf. 720) shows that $\unsure{S}$ -- over \uline{any} \uuline{base}
\diagram{\begin{tikzcd}[ampersand replacement=\&,column sep=small]
H^0(S,\mathfrak S_3,\mathcal J^0)\arrow[d,equal]\arrow[r,phantom,"\simeq"] \& H^1(\deleted{\unsure{S,}}\mathfrak S_3;H^0(\unsure{S|S_0},\mathcal J^0))\simeq0\quad\unsure{\Leftarrow}\\
H^0(\mathfrak S_3,H^0(S,\mathcal J^0))
\end{tikzcd}}}

(\unsure{the} only cube \deleted{\unsure{of}} root of 1 in $\mathbb Z$ (or even in $\mathbb Q$) is 1. This establishes the first relation (307). For the second, we use the spectral sequence

\margin{146}

% Archive printed footer (not manuscript):
% 146
% https://grothendieck.umontpellier.fr
% Université de Montpellier, UMR CNRS 5149, Case courrier 051, Place Eugène Bataillon
% 34095 Montpellier cedex 5 - France

\scanpage{148}

\diagram{\begin{tikzcd}[ampersand replacement=\&]
H^{\unsure{*|n}}(S_0,\mathfrak S_3;\mathcal J^0) \&
\overset{\inserted{E_2^{pq}=}}{H^p(\mathfrak S_3,\uline{H}^{q}(S_0,\mathcal J^0))}\arrow[l,Rightarrow]
\end{tikzcd}}
which gives the exact sequence
\diagram{\begin{tikzcd}[ampersand replacement=\&,column sep=small]
0\arrow[r] \& H^1(\mathfrak S_3,H^0(S_0,\mathcal J^0))\arrow[r,"\unsure{\deleted{\simeq}|\simeq}"] \& H^1(S_0,\mathfrak S_3;\mathcal J^0)\arrow[r] \& H^0(\mathfrak S_3,H^1(S_0,\mathcal J^0))\arrow[d]\arrow[dl,equal]\\
\& \& \bigl(H^{\unsure{0|q}}(\mathfrak S_3,M_{\unsure{0}}\otimes_{\mathbb Z}H^1(S_0,\mathbb G_m))\bigr) \& H^2(\mathfrak S_3,H^0(S_0,\mathcal J^0))
\end{tikzcd}}

Now we have
\[
H^1(S_0,\mathcal J^0)\simeq H^1(S_0,\mathbb G_m^2)\simeq\operatorname{Pic}(\unsure{S|S_0})^2=0
\]
\margin{\uline{N.B.} we do \uline{not} need the relation $\operatorname{Pic}(S)=0$}
hence we get
\[
H^1(S_0,\mathfrak S_3;\mathcal J^0)\simeq H^1(\mathfrak S_3,\underbrace{H^0(S_0,\mathcal J^0)}_{\substack{\deleted{(\mathbb Z^*)^{\{0,1,\infty\}}/\mathbb Z^*}\\\deleted{\simeq\mathbb Z^*\otimes_{\mathbb Z}M_0}}})
\]

But in fact, the calculations made elsewhere show that we have
\[
\tag{309}
H^0(\mathfrak S_3,M_{\unsure{0|\alpha}}\otimes_{\mathbb Z}\unsure{I|I_0})=H^1(\mathfrak S_3,M_{\unsure{0|\alpha}}\otimes_{\mathbb Z}I)=0
\]
for any $\mathbb Z$-module $I$ (for example $I=\unsure{J^0|\mathcal J^0}(S)$, $S$ an \unsure{arbitrary|arbitrary} scheme, \inserted{or $I=H^1(S,\mathbb G_m)=\operatorname{Pic}(S)$}), whence finally,
\[
\tag{\unsure{310}}
H^0(\unsure{\deleted{S_0}|S},\mathfrak S_3;\mathcal J^0)\simeq H^1(S,\mathfrak S_3,\mathcal J^0)=0
\]
for $\forall$ scheme $S$. We thus find:

\margin{147}

% Archive printed footer (not manuscript):
% 147
% https://grothendieck.umontpellier.fr
% Université de Montpellier, UMR CNRS 5149, Case courrier 051, Place Eugène Bataillon
% 34095 Montpellier cedex 5 - France

\scanpage{149}
\mspage{722}

\uline{Proposition} Let \deleted{$\unsure{\mathcal J}$}
$\mathcal J^0=\mathbb G_m^{\{0,1,\infty\}}/\underbrace{\mathbb G_m}_{\mathrm{diag.}}\simeq\mathbb G_m\otimes_{\mathbb Z}\unsure{M_0|M}$
Then for $\forall$ $(\mathfrak S_3,\mathcal J^0)$-torsor over an arbitrary base $S$, there exists a unique trivialization of it, i.e. a unique section, invariant under the operations of $\mathfrak S_3$.

\uline{Corollary 1} Over any base $S$, there exists a unique morphism
\[
(311)\qquad
\diagram{\begin{tikzcd}[ampersand replacement=\&]
\psi:U_{0,3\,S} \arrow[r] \& \mathcal J^0_S
\end{tikzcd}}
\]
which is compatible with the operations of $\mathfrak S_3$, and which (via the homomorphism
$\mathcal J^*_S\longrightarrow\mathcal J^0_S$ \unsure{which|that} \unsure{it|it} \unsure{is factored by}, by the universal property of the \illegible{1} generalized \unsure{Jacobians}) induces the identity homomorphism
$\underset{\inserted{\mathcal J^*_S}}{\mathcal J^0_S}-\mathcal J^0_S$.

\margin{\uline{N.B.} With the notation of \unsure{695}, we have $U_{0,3\,S}\simeq(U_1)_S$, where $U_1\subset\mathbb G_m^3/\mathbb G_m$ is \unsure{defined} by
\[
(312)\qquad
U_1=\{(x_0,x_1,x_\infty)\in\Gamma\mathbb G_m^{\{0,1,\infty\}}\mid x_0+x_1+x_\infty=\unsure{0}\}
\quad\mathrm{mod}\ \mathrm{diag}_3\mathbb G_m.
\]}

\uline{Corollary 3} Let $Y$ be a scheme of

\margin{148}

% Archive printed footer (not manuscript):
% 148
% https://grothendieck.umontpellier.fr
% Université de Montpellier, UMR CNRS 5149, Case courrier 051, Place Eugène Bataillon
% 34095 Montpellier cedex 5 - France

\scanpage{150}

\unsure{Brauer}-Severi over $S$, of relative dimension $1$ (i.e. a form of $\mathbb P^1_S$), equipped with a \inserted{closed} subscheme $S'\subset Y$, étale of degree $3$ over $S$, and let \deleted{$\unsure{X}$} $U=Y\setminus S'$, $\mathcal J^*_{U/S}$ the generalized Jacobian, so that we have
\[
(313)\qquad
\diagram{\begin{tikzcd}[ampersand replacement=\&,column sep=1.5em,row sep=.5em]
1 \arrow[r,no head] \& \mathcal J^0_{U/S} \arrow[r,no head] \arrow[d,phantom,"\unsure{\wr}" description] \& \mathcal J^*_{U/S} \arrow[r,no head] \& \mathbb Z_S \arrow[r,no head] \& 1 \\
\& \mathbb G_m\otimes_{\mathbb Z}\unsure{M_0}(S'/S) \& \& \&
\end{tikzcd}}
\]
where $\unsure{M_0}(S'/S)\simeq\unsure{p_*|\rho_*}(\mathbb Z_{S'})/\underbrace{\mathbb Z_S}_{\mathrm{diag}}$ is the $S$-scheme \inserted{of \deleted{groups} \inserted{$\mathbb Z$-modules}} (locally $\unsure{\overline{\mathrm{isom.}}}$ to $\mathbb Z^{\unsure{2}}$, deduced from $\unsure{M_{0\,S}}$ by twisting it (via the action of $\mathfrak S_3$) by the \deleted{principal bundle} \inserted{$\mathfrak S_3$-torsor over $S$} associated with $S'$.

There exists then a unique morphism
\[
(\unsure{314})\qquad
\diagram{\begin{tikzcd}[ampersand replacement=\&]
\psi:U \arrow[r] \& \mathcal J^0_{U/S}
\end{tikzcd}}
\]
commuting with \deleted{$\uline{\operatorname{Aut}}$} the action of $\uline{\operatorname{Aut}}(S'/S)\xleftarrow{\sim}\uline{\operatorname{Aut}}_S(Y,S')$, and inducing (by the universal property) \deleted{\unsure{a}} the \unsure{composite} homomorphism
\[
\diagram{\begin{tikzcd}[ampersand replacement=\&]
\mathcal J^0_{U/S} \arrow[r,hook] \& \mathcal J^*_{U/S} \arrow[r,"\widetilde\psi"] \& \mathcal J^0_{U/S}
\end{tikzcd}}
\]
which is the identity.

\margin{149}

% Archive printed footer (not manuscript):
% 149
% https://grothendieck.umontpellier.fr
% Université de Montpellier, UMR CNRS 5149, Case courrier 051, Place Eugène Bataillon
% 34095 Montpellier cedex 5 - France

\scanpage{151}
\mspage{723}

Let us make explicit \unsure{the|the} \deleted{homo}morphism
\[
(315)\qquad
\diagram{\begin{tikzcd}[ampersand replacement=\&,column sep=2em,row sep=.9em]
U_{0,3}=\mathbb P^1\setminus\{0,1,\infty\}_{S_0}
  \arrow[r,hook,"\psi"] \arrow[d,equal]
\& \underbrace{\mathbb G_m^{\{0,1,\infty\}}/\mathbb G_m}_{\unsure{\mathcal J^0}}
\\
\begin{gathered}
\deleted{\{z\in\unsure{\Gamma}\mathbb G_m\mid z-1\ \mathrm{inv}\}}\\
\unsure{(}\mathbb G_m\setminus\{1\}\unsure{)}_{S_0}
\simeq\{z\in\Gamma\mathbb G_m\mid z-1\ \mathrm{inv}\}
\end{gathered}
\& \begin{gathered}
\deleted{\unsure{\mathcal J^1}}\inserted{U_1=\{(x_0,x_1,x_\infty)^{\deleted{\unsure{\mathbb Z}}}
\in\Gamma\mathbb G_m^3/\mathbb G_m\mid}\\
\inserted{x_0+x_1+x_\infty=0\}}
\end{gathered} \arrow[u,hook]
\end{tikzcd}}
\]
\deleted{\unsure{it}} \deleted{\unsure{it is}} isomorphism of $U_{0,3}$ with $U_1$ which \deleted{\unsure{is}} corresponds: \deleted{\unsure{one must|extends}}
\[
(316)\qquad
\left\{\begin{array}{l}
\diagram{\begin{tikzcd}[ampersand replacement=\&,column sep=2em,row sep=.5em]
\overline\psi:\unsure{\widehat U_{0,3}|\overline U_{0,3}} \arrow[r] \arrow[d,equal]
\& \widehat U_1\subset\mathbb P^2_{S_0} \arrow[d,equal]\\
\mathbb P^1_{S_0}
\& \{(x_0,x_1,x_\infty)\in\unsure{\widehat{\mathbb P}^2_{\unsure{S_0}}|\mathbb P^2_{S_0}}\mid x_0+x_1+x_\infty=0\}
\end{tikzcd}}\\[1em]
\overline\psi(0)=(0,1,-1)\\
\overline\psi(1)=(-1,0,1)\\
\overline\psi(\infty)=(1,-1,0)
\end{array}\right.\ ,
\]
it suffices to take
\[
(317)\qquad
\psi(z)=(-z,z-1,1)=(1,1/z-1,-1/z)
\inserted{=\left(\frac{-z}{z-1},1,\frac{1}{z-1}\right)}.
\]

Let us now return to the datum of a relative elliptic curve \inserted{$X$} over $S$ modulo symmetry (corresponding to a morphism
$S\longrightarrow\deleted{\unsure{\operatorname{Spec}}}\mathcal M'_{1,1}$).

\margin{150}

% Archive printed footer (not manuscript):
% 150
% https://grothendieck.umontpellier.fr
% Université de Montpellier, UMR CNRS 5149, Case courrier 051, Place Eugène Bataillon
% 34095 Montpellier cedex 5 - France

\scanpage{152}

We assume $2$ invertible on $S$, so that the datum of $X$ \inserted{modulo $\omega_X$} corresponds, by the Legendre \unsure{carpet} $X\mapsto(X/\deleted{\text{\unsure{modulo symmetry}}}\inserted{\{1,\omega_X\}})$, to the datum of a $(Y,S')$ as above, with in \uline{addition} a section of $U=Y-S'$.

In fact, the datum of $(Y,S')$ amounts to the datum of $S'$, \inserted{i.e. of an $\mathfrak S_3$-torsor} \deleted{$Y$ is recovered} $T$ over $S$ --- and $Y$ is recovered from \unsure{there|it} as the twisted scheme $T\wedge_{\mathfrak S_3}\mathbb P^1_S$ \deleted{$\unsure{\wedge_{\mathfrak S_3}}$ \unsure{and|or}, $T\wedge_{\mathfrak S_3}\mathcal J^0_{(U_{0,3})_S}\simeq\unsure{\mathcal J^0}$}
with
\[
U=Y\setminus S'\simeq(T\wedge_{\mathfrak S_3}\mathbb P^1_S)\setminus\underbrace{T\wedge_{\mathfrak S_3}\{0,1,\infty\}_S}_{S'}
\]
canonically embedded in
\[
T\wedge_{\mathfrak S_3}\mathcal J^0_{(U_{0,3})_S}\simeq\mathcal J^0_{U/S}.
\]
The embedding $U\subset\mathcal J^0_{U/S}$ can also be interpreted as a canonical morphism
\[
(318)\qquad
\diagram{\begin{tikzcd}[ampersand replacement=\&]
X \arrow[r] \& \mathcal J^0_{U/S}\underbrace{=\unsure{\mathcal J^0_{0\,X/S}|\mathcal J^0_{X/S}}}_{\substack{\text{\unsure{with|in} notation}\\\text{of page }720}}
\end{tikzcd}}
\]
which should be easy to characterize \unsure{axiomatically}

\margin{151}

% Archive printed footer (not manuscript):
% 151
% https://grothendieck.umontpellier.fr
% Université de Montpellier, UMR CNRS 5149, Case courrier 051, Place Eugène Bataillon
% 34095 Montpellier cedex 5 - France